\section{Motivation and Challenges}
\label{sec:motivation}

The resulting management problem is to decide which KV to retain and
in which tier. For a multi-agent service, the benefit of these decisions
depends on how retained state supports both normal execution and workflow
completion after a failure. Existing fault-tolerant serving and
workflow-aware KV management provide useful mechanisms, but their
objectives leave this joint decision unresolved.

\subsection{Limitations of Existing Approaches}
\label{sec:motivation-limitations}

\begin{figure*}[t]
  \centering
  \captionsetup[subfigure]{font=normalsize,labelfont=normalfont,
    textfont=normalfont,justification=centering,singlelinecheck=false}
  \begin{subfigure}[t]{2.48in}
    \vspace{0pt}
    \begin{minipage}[c][1.65in][c]{\linewidth}
      \centering
      \includegraphics[width=\linewidth]{figures/workflow.pdf}
    \end{minipage}
    \caption{Four-branch workflow}
    \label{fig:motivation-workflow}
  \end{subfigure}\hfill
  \begin{subfigure}[t]{1.86in}
    \vspace{0pt}
    \begin{minipage}[c][1.65in][c]{\linewidth}
      \centering
      \includegraphics[width=\linewidth]{figures/workflow-gaps.pdf}
    \end{minipage}
    \caption{Workflow-maximum gaps}
    \label{fig:motivation-branch-gap}
  \end{subfigure}\hfill
  \begin{subfigure}[t]{2.52in}
    \vspace{0pt}
    \begin{minipage}[c][1.65in][c]{\linewidth}
      \centering
      \includegraphics[width=\linewidth]{figures/recovery-exposure.pdf}
    \end{minipage}
    \caption{Recovery cost and WCT}
    \label{fig:motivation-recovery-exposure}
  \end{subfigure}
  \caption{Workflow timing and conditional recovery impact.
    (a) Four trace-derived parallel agents A1-A4 with schematic Main agents A0/A5.
    Bars preserve timing; dashed tails assume an all-branch join.
    (b) Empirical CDF of workflow-level gaps over 121 workflows. For each parallel group with temporally inferred fan-in, we subtract the shortest agent execution time from the longest, then retain the largest gap per workflow. 
   (c) Each horizontal bar represents one agent (A1--A4) and shows its
extra delay to reconstruct KV after a failure, measured by full
prefill with Qwen2.5-7B-Instruct-1M. The hatched portion is absorbed
by slack; the solid portion is the calculated WCT increase.}

  \label{fig:motivation}
\end{figure*}

\paragraph{Limitation 1: Fault tolerance without workflow-aware priorities.}
Fault-tolerant serving systems preserve state and coordinate recovery
for interrupted requests. D\'ej\`aVu streams KV checkpoints, and GhostServe
uses host-memory parity to reduce reconstruction
work~\cite{strati2024dejavu,jayakody2026ghostserve}.
LUMEN further coordinates checkpoint placement, recovery routing and
capacity restoration according to worker load~\cite{cao2026lumen}.
These mechanisms address state loss and recovery contention, but do not
use multi-agent dependencies to determine which requests' faster
recovery would advance workflow completion. Request recovery cost and
worker load alone cannot express this benefit.

\label{sec:motivation-workflow}
Figure~\ref{fig:motivation}(\subref{fig:motivation-workflow}) illustrates
why workflow position matters. A2 finishes last, 151.6\,s after the fork,
whereas A1 finishes at 60.4\,s. If the following Main stage requires all
four outputs, A1 can absorb 91.2\,s of additional delay before postponing
the join, whereas any delay on A2 postpones it. This is an example from the Claude code traces~\cite{semianalysis2026cctraces}, which contains the workflows that are very similar to the architecture shown in part a of \ref{fig:agent-workflow-examples}. We call the additional
delay a branch can absorb without postponing the required output its
\emph{slack}. 

This situation where agents with parallel execution have different completion times is very common in multi-agent applications, and the time differences can be significant. We analyze the Claude code trace~\cite{semianalysis2026cctraces} related information and visualize this time difference in Figure~\ref{fig:motivation}(\subref{fig:motivation-branch-gap}).
This figure reports each workflow's maximum longest-minus-shortest branch-duration gap,
with a median of 138\,s and a 90th percentile of 1,173\,s. When the Main agent waits for all required branches, differences in branch completion times can leave shorter branches with slack. Additional delays within this slack do not postpone the join or
increase WCT.

These differences in slack affect how KV reconstruction delays workflow
completion after a failure.
Figure~\ref{fig:motivation}(\subref{fig:motivation-recovery-exposure})
shows reconstruction times measured by full prefill with
Qwen2.5-7B-Instruct-1M using sparse
attention~\cite{qwen2025qwen25oneM,jiang2024minference},
and their calculated impact on WCT.
We consider each agent's interruption separately, with other execution
unchanged and the Main agent waiting for all four branches.
A1's 33.2\,s reconstruction delay is fully absorbed by its 91.2\,s slack,
leaving WCT unchanged. In contrast, A2 has no slack, so its 32.6\,s
reconstruction delay increases WCT by the same amount.
Thus, similar reconstruction costs can have very different effects
on workflow completion. Recovery priorities should therefore account
for a request's position in the workflow, in addition to its
reconstruction cost.

\paragraph{Limitation 2: Workflow-aware caching without fault tolerance.}
Workflow-aware systems already exploit the dependencies highlighted
above to improve application execution. Parrot exposes inter-request
dependencies for application-level
scheduling~\cite{lin2024parrot}. KVFlow guides prefix retention and
prefetching by upcoming agent invocations~\cite{pan2025kvflow}, while Pythia
predicts future workflow steps and prompt contents to manage and stage
KV~\cite{yu2026pythia}. TOPAS jointly selects resident prefixes and admitted
requests by balancing remaining-path progress, reuse, and movement or
preemption costs~\cite{ni2026topas}. These workflow-aware caching and
scheduling policies do not address serving-worker failures and the
resulting loss of in-progress KV.

Worker failures must be accounted for in large-scale deployment.
As discussed in Section~\ref{sec:motivation-background}, worker failures
can make KV unavailable and interrupt several co-located
agents~\cite{strati2024dejavu,jayakody2026ghostserve,cao2026lumen}.
For an agent that has accumulated context over many requests, resuming
execution can require rebuilding a long history.
Figure~\ref{fig:motivation-recompute} shows that this reconstruction can
take far longer than reusing retained KV. The preceding workflow example
shows how the resulting delay can postpone required outputs.
Failure handling is therefore necessary for workflow execution,
even when normal-operation caching is effective.

Offloading KV provides useful building blocks for this protection.
CachedAttention and Mooncake extend retention beyond GPU
memory~\cite{gao2024cachedattention,qin2025mooncake}; Pythia writes
future-relevant contexts to shared storage accessible by other
workers~\cite{yu2026pythia}. A surviving copy can support recovery, but
retaining contexts for future invocations does not ensure that an
interrupted agent's state has been protected. For example, KVFlow
protects fixed prompts while assigning varying suffixes the highest
eviction priority~\cite{pan2025kvflow}. Those suffixes can still be needed
to resume the interrupted request, even if later requests would not
reuse them. Shared prefixes may therefore survive while the agent's
task-specific history still needs reconstruction.

\label{sec:motivation-long-context}
The two limitations thus expose complementary gaps: existing
fault-tolerant serving does not incorporate multi-agent dependencies
into recovery priorities, while the workflow-aware caching policies
above do not incorporate worker failures.
A multi-agent service needs KV management that supports both normal
workflow execution and recovery from state loss. This requires
combining workflow information with fault-tolerant state protection
under finite storage and bandwidth budgets.

\subsection{Design Challenges}
\label{sec:motivation-challenges}

Together, these limitations call for a policy that connects workflow
progress, failure recovery and KV storage management. Closing the gap
requires estimating the value of retained state, translating that value
into placement decisions, and maintaining those decisions as execution
changes. These are coupled challenges: the usefulness of fast recovery
depends on workflow progress, while protection itself consumes resources
needed by the workflow.

\paragraph{Challenge 1: Estimating the value of KV protection online.}
A protection decision must relate saved reconstruction or retrieval time
to the workflow's required outputs. The example above obtains slack from
completed traces; an online system must act before the remaining model
and tool durations are known. The branch that gates completion can
change, and one worker failure can interrupt several agents whose
recoveries share resources. Normal reuse also contributes value even
when recovery is not urgent. The challenge is to derive useful priorities
from partial execution and reuse history without treating request-level
savings as independent, additive WCT reductions.

\paragraph{Challenge 2: Translating protection value into tiered placement.}
A request's urgency alone does not specify how much of its KV should
occupy each tier. A policy must distinguish retaining enough state to
avoid reconstruction from using faster capacity to reduce retrieval
latency. It must compare these benefits with the bytes retained, the
capacity cost of each tier and the work needed to move state.
The same KV may support several agents, while different portions of a
long history may have different access demands. Allocating capacity
therefore requires connecting workflow-level priorities to the state
that recovery and reuse actually need, subject to finite tier budgets.

\paragraph{Challenge 3: Adapting protection without disrupting inference.}
Agent progress, tool returns and new requests change both protection
value and foreground memory demand. Repeatedly reassessing every stored
KV object can make fine-grained management expensive; frequent placement
changes can also consume bandwidth needed by active requests.
A useful policy must update protection with bounded decision and
migration work, react to capacity pressure, and retain valid recovery
paths while state moves between tiers. The challenge is to keep
protection responsive without allowing its management and data movement
to erase the workflow benefit it seeks.

\label{sec:motivation-goals}
These challenges motivate AegisKV, a workflow-aware KV protection system
that coordinates recovery priorities, tiered placement and online state
management. The next section develops the observations and mechanisms
behind this design.
